\section*{4. de Rham 上同调}

\subsubsection*{4.1.}
设 $D$ 是复解析圆盘，$0\in D$，令 $D^*=D-\{0\}$，并设 $\widetilde D^*$ 是 $D^*$ 的一个万有覆盖。设 $X$ 是解析空间，$f:X\longrightarrow D$ 是态射，并考虑 (1.2.7.1) 的图：
\begin{equation}\tag{4.1.1}
\begin{tikzcd}[column sep=large,row sep=large]
X_0 \arrow[r,hook,"\bar i"] \arrow[d,equal]
  & \bar X \arrow[d,"p"']
  & \bar X^* \arrow[l,hook',"\bar j"'] \arrow[d,"p'"'] \\
X_0 \arrow[r,hook,"i"]
  & X
  & X^* \arrow[l,hook',"j"']
\end{tikzcd}
\end{equation}
假设 $X^*=X-X_0$ 非奇异（即光滑）。因此其覆盖 $\bar X^*$ 也光滑。用 $\Omega_{\bar X^*}^*$ 表示 $\bar X^*$ 上全纯微分形式层的 de Rham 复形。它是常值层 $\mathbb C$ 的消解（全纯 Poincaré 引理），因而给出 $D^+(\bar X,\mathbb C)$ 中的态射
\begin{equation}\tag{4.1.2}
\bar j_*\Omega_{\bar X^*}^*
\longrightarrow R\bar j_*\mathbb C
\end{equation}

\subsubsection*{命题 4.2.}
\underline{态射 (4.1.2) 是同构。}

显然，(4.1.2) 在 $\bar X^*\subset\bar X$ 的每一点处都是拟同构。\footnote{原印本的证明写作 (4.1.1)，那是前面的图而非所讨论的态射；英文恢复了显然应有的引用 (4.1.2)。}
因此，只需证明每个点 $x\in X_0$ 在 $\bar X$ 中都有一个基本邻域系 $U$，使
\[
H^q\bigl(\bar j^{-1}(U),\Omega^p\bigr)=0
\qquad\text{当 }q>0.
\]
当 $V$ 遍历 $X$ 中 $x$ 的一个基本邻域系时，各集合 $p^{-1}(V)$ 构成 $\bar X$ 中 $x$ 的一个基本邻域系。特别地，可取 $p^{-1}(V)$，其中 $V$ 是 $X$ 中 $x$ 的一个 Stein 邻域。若 $V$ 为 Stein，则 $V-X_0$ 仍为 Stein（因为 $X_0\cap V$ 在 $V$ 中由一个方程定义），其覆盖
\[
p'^{-1}(V-X_0)=\bar j^{-1}p^{-1}(V).
\]
也为 Stein。因此，对 $\bar X^*$ 上的每个凝聚层 $\mathcal F$，
\[
H^q\bigl(\bar j^{-1}p^{-1}(V),\mathcal F\bigr)=0
\qquad\text{当 }q>0,
\]
命题得证。

\subsubsection*{推论 4.3.}
\underline{同构 (4.1.2) 诱导同构}
\begin{equation}\tag{4.3.1}
\bar i^*\bar j_*\Omega_{\bar X^*}^*
\xrightarrow{\ \sim\ }R\Psi_\eta(\mathbb C).
\end{equation}

\subsubsection*{变体 4.4.}
\underline{设 $V$ 是 $X^*$ 上的有限维复向量空间局部常值层（复局部系统）。Poincaré 引理的拟同构}
\[
V\longrightarrow\Omega_{X^*}^*(V)
\]
\underline{诱导拟同构}
\begin{equation}\tag{4.4.1}
\bar i^*\bar j_*\Omega_{\bar X^*}^*(p'^*V)
\xrightarrow{\ \sim\ }R\Psi_\eta(V).
\end{equation}

我们将证明：将 $\bar i^*\bar j_*\Omega_{\bar X^*}^*$ 换为由增长条件和单值条件定义的一个子复形后，(4.3.1) 仍成立。